\section{Metoodika} \label{metoodika}

Käesolev peatükk kirjeldab, kuidas töö eesmärki realiseeritakse. Alapeatükis~\ref{metoodika:protsess} esitatakse uurimisprotsess tervikuna ning eristatakse selgelt, millised sammud on käesoleva töö raames lõpule viidud ja millised on kavandatud jätkukavasse. Alapeatükk~\ref{metoodika:kategooriad} määratleb tagasisidekategooriad operatsionaalselt. Alapeatükis~\ref{metoodika:mudel} põhjendatakse mudeli valikut. Alapeatükk~\ref{metoodika:promptid} esitab kahe võrreldava prompti disaini. Alapeatükis~\ref{metoodika:hindamine} kirjeldatakse \emph{kavandatud} täismahulise hindamise metoodikat ja testkogu disaini.

\subsection{Uurimisprotsess} \label{metoodika:protsess}

Uurimine järgib disain\-teaduse (\emph{ingl. design science research}) põhimõtet~\cite{hevner_design_2004}: identifitseeritakse probleem, projekteeritakse artefakt (käesoleval juhul prototüüp), rakendatakse seda ning hinnatakse rangelt. Hindamise põhjal saab teha järeldusi nii artefakti kui ka selle aluseks olevate disainivalikute kohta.

Töö protsess on kavandatud koosnema järgmistest etappidest:
\begin{enumerate}
    \item Tagasisidekategooriate operatsionaalne määratlemine kirjanduse ja TÜ juhendi põhjal.
    \item Üldise ja struktureeritud prompti disain.
    \item Prototüübi realiseerimine veebirakendusena, mis kutsub välja LLM-i API-d ning toetab salvestatud vastustega demo-režiimi.
    \item Testkogu disain: 20 eestikeelset informaatika lõputöö katkendit erinevatest peatükitüüpidest, koos kuldstandardi koostamise protseduuriga.
    \item Pilootuuring väikesel tekstivalimil: prototüübi väljundi kvalitatiivne analüüs prompti disaini iteratsiooniks.
    \item Testkogu täielik annoteerimine kuldstandardiks, sealhulgas teise sõltumatu annoteerija kaasamine ja annoteerijatevahelise kokkuleppe arvutamine.
    \item Prototüübi täismahuline käivitamine kogu testkogu vastu mõlema prompti versiooniga ja mõlema mudeliga (Claude 4.7 Opus, GPT-5.5); mudeli väljundi võrdlemine kuldstandardiga ja näitajate (täpsus, saagis, F\textsubscript{1}) arvutamine iga kategooria kohta eraldi.
    \item Üldise ja struktureeritud prompti tulemuste kvantitatiivne võrdlus.
\end{enumerate}

\textbf{Käesoleva töö ulatus.} Sammud 1--5 on \emph{käesoleva töö raames lõpule viidud}. Sammud 6--8 jäävad teadlikult käesoleva töö ulatusest välja ning on kavandatud jätkukavasse (vt alapeatükk~\ref{hindamine:edasine}). See piiritlemine on tehtud kahel põhjusel. Esiteks võimaldab see käesoleva töö raames keskenduda kvalitatiivselt prototüübi väljundi mustrite ja prompti disaini iteratsiooni põhjalikule analüüsile, mis on edasise kvantitatiivse hindamise jaoks vajalik eeldus. Teiseks ei oleks sõltumatu teise annoteerija kaasamine ning 20-katkendiline annoteerimine käesoleva töö esitamise ajaraamis piisava kvaliteediga teostatav, ning ebakvaliteetne annoteerimine annaks ebakvaliteetse kvantitatiivse tulemuse, mis kahjustaks tervet uuringut.

\subsection{Tagasisidekategooriate operatsionaalne määratlus} \label{metoodika:kategooriad}

Töös käsitletakse nelja tagasisidekategooriat. Iga kategooria kohta on tabelis~\ref{tab:kategooriad} esitatud määratlus, näide tüüpilisest probleemist ning probleemi puudumise näide. Need määratlused olid aluseks nii prompti koostamisele kui ka tulevase kuldstandardi annoteerimisreeglitele.

\begin{table}[htb!]
    \centering
    \caption{Tagasisidekategooriate operatsionaalsed määratlused.}
    \label{tab:kategooriad}
    \begin{tblr}{width=1.0\textwidth, hlines, vlines,
                colspec = { Q[l,font=\bfseries,wd=0.18\textwidth] X[l] X[l] },
                row{1} = {font=\bfseries, bg = colorCellHighlight},
            }
    Kategooria & Tüüpiline probleem & Probleemi puudumise tunnus \\
    Struktuur & Peatükk algab kohe alapealkirjaga ilma sissejuhatava lõiguta; lõikude vahel ei ole loogilist üleminekut. & Iga peatükk algab sissejuhatava lõiguga; alampeatükid on omavahel sujuvalt seotud. \\
    Akadeemiline stiil & Kõnekeel (\emph{tüüp, asi}); liiga pikad laused; mina-vormi liigne kasutamine; otsetõlke ilmingud. & Kasutatakse neutraalset akadeemilist stiili; lauseehitus on selge. \\
    Terminoloogia & Sama mõiste tähistamiseks kasutatakse läbisegi mitut eri terminit (\emph{lõim} ja \emph{niit}); ingliskeelseid termineid esitatakse ilma esmakasutuse selgituseta. & Terminid on järjepidevad; ingliskeelsed terminid esitatakse esmakasutusel ka eestikeelselt. \\
    Viitamisvajadus & Konkreetsed arvulised või faktilised väited on esitatud ilma viiteta; metoodika kirjeldus on esitatud kui üldteadmine. & Iga väide, mis pole autori enda originaalne mõte, on viidatud. \\
    \end{tblr}
\end{table}

Kategooriad on disainitud nii, et nad on \emph{vastastikku eristatavad} -- ühte sõnastust ei tohiks korraga liigitada kahte kategooriasse. Seda eristatavust kontrolliti pilootuuringus mudeli väljundi käsitsi läbivaatamise käigus ning eristatavus osutus piisavaks; täismahulises annoteerimises kontrollitakse seda annoteerijatevahelise kokkuleppe abil.

\textbf{Pilootuuringust tulenevad määratluste täpsustused.} Pilootuuringus täheldatud näilised valepositiivsed (vt alapeatükk~\ref{hindamine:kategooriad}) viisid kahe määratluse selgesõnalise täpsustuseni, mida rakendatakse nii prompti vastunäidetes kui ka kuldstandardi annoteerimisreeglites:
\begin{itemize}
    \item \emph{Viitamisvajadus} ei rakendu üldteadmise tasandi väidetele -- näiteks levinud protokollide või põhimõistete tutvustamisele --, mille puhul viidet ei oleks ka inimese kirjutatud akadeemilises tekstis oodatud.
    \item \emph{Akadeemiline stiil} ei märgi probleemiks mina-vormi kasutust kohtades, kus tudeng kirjeldab selgelt enda originaalset panust või metoodilist valikut, kuna Tartu Ülikooli juhend lubab mina-vormi sellises kontekstis.
\end{itemize}
Need täpsustused on \emph{tulemus} pilootuuringust ning on aluseks täismahulise hindamise annoteerimisreeglitele.

\subsection{Mudeli valik} \label{metoodika:mudel}

Töös kasutatakse esmase mudelina Anthropici \texttt{claude-opus-4-7} mudelit. Valiku tegemise kriteeriumid olid:
\begin{itemize}
    \item \textbf{Eesti keele tugi}: TartuNLP avaliku eesti keele LLM-võrdlusplatvormi (baromeeter)~\cite{tartunlp_baromeeter_2025} põhjal kuuluvad tipptasemel Claude'i ja GPT mudelid eesti keele kvaliteedi poolest tippmudelite hulka.
    \item \textbf{Pikk konteksti aken}: 1\,000\,000 tokenit~\cite{anthropic_models_2026} võimaldab analüüsida kogu lõputöö peatükki korraga ilma sisendi tükeldamiseta.
    \item \textbf{Struktureeritud väljundi tugi}: mudel järgib usaldusväärselt JSON-skeemi kohaseid väljundinõudeid (vt alapeatükk~\ref{prototuup:promptid}).
    \item \textbf{Hind ja kättesaadavus}: API-ligipääs on tudengieelarve piires.
\end{itemize}

Tulemuste üldistatavuse kontrollimiseks katsetatakse paralleelselt ka OpenAI \texttt{gpt-5.5} mudelit, mille tugi on prototüübis realiseeritud. Alapeatükis~\ref{hindamine:synteetiline} esitatud empiirilise hindamise raames käivitati prototüüp kõikide 20 sünteetilise testkatkendi vastu mõlema mudeli ja mõlema prompti versiooniga (kokku 80 API-päringut). Mõlema uue põlvkonna mudeli (Claude 4.7 Opus ja GPT-5.5) Messages/Chat Completions API ei aktsepteeri enam \texttt{temperature}-parameetrit, mistõttu pakkujakihid jätavad selle parameetri päringust automaatselt välja ning kasutatakse mudelite vaikeväärtuseid~\cite{anthropic_models_2026, openai_gpt5_2026}.

\subsection{Promptide disain} \label{metoodika:promptid}

Töös võrreldakse kahte prompti varianti: \emph{üldist} ja \emph{struktureeritud}. Mõlemad nõuavad mudelilt sama JSON-skeemiga vastust (et automaatne järeltöötlus oleks võrreldav), kuid erinevad mahukalt selles, kui palju juhiseid mudelile kategooriate ja oodatava käitumise kohta antakse.

\subsubsection{Üldine prompt}

Üldine prompt on lühike: see annab mudelile rolli ja JSON-skeemi (sealhulgas vabavormilise \texttt{MUU}-kategooria), kuid \emph{ei sisalda} ühegi kategooria definitsiooni, näiteid, vastunäiteid ega kindlushinnangu täpsemat juhist. Seega on üldine prompt funktsionaalselt minimaalne tasapind, millele struktureeritud prompt lisab kategooriapõhise sisu:

\begin{minted}[breaklines,fontsize=\small]{text}
Oled eestikeelse informaatika lõputöö juhendaja abiline. Anna sisestatud
peatükile tagasisidet selle kvaliteedi kohta. Ole konstruktiivne ja
konkreetne. Tagasiside peab olema eesti keeles.

OLULINE: Sa ei kirjuta lõputööd tudengi eest ega sõnasta lõike ümber.
Anna ainult soovituslikku tagasisidet.

Vasta JSON-formaadis järgmise skeemi järgi:
{
  "leiud": [
    {
      "kategooria": "STRUKTUUR" | "AKADEEMILINE_STIIL"
                  | "TERMINOLOOGIA" | "VIITAMISVAJADUS" | "MUU",
      "tsitaat":  "tekstis täpne probleemne kohaviit (kuni 200 tähemärki)",
      "probleem": "lühike kirjeldus, mis selles kohaviites on valesti",
      "põhjendus": "miks see on probleem",
      "soovitus": "konkreetne soovitus (kuni 200 tähemärki)",
      "kindlus":  "kõrge" | "keskmine" | "madal"
    }
  ]
}

Peatüki tüüp: ${peatuki_tyyp}
Peatüki tekst:
"""
${tekst}
"""

Vasta ainult JSON-iga, ilma lisaselgituseta.
\end{minted}

\subsubsection{Struktureeritud prompt}

Struktureeritud prompt erineb üldisest neljal olulisel viisil: (1) iga tagasisidekategooria täpne määratlus koos vähehetk-näidetega ja vastunäidetega, (2) selgesõnaline eetiline meelespea (mudel ei tohi ümber sõnastada), (3) mõtteahel-element (\emph{põhjendus}-väli, mis nõuab mudelilt enda järelduse põhjendamist enne soovitust) ja (4) reeglid, mis suunavad mudelit konservatiivsusele (\emph{parem vähem õigeid kui palju valesid}). Kategooria \texttt{MUU} on jäetud välja, nõudes mudelilt nelja täpselt määratletud kategooria kasutamist. Prompti täielik tekst on esitatud lisas~II.

Prompti võtmeosad on järgmised:

\begin{minted}[breaklines,fontsize=\small]{text}
Oled eestikeelse informaatika lõputöö tagasisidesüsteem.
SUL EI OLE LUBATUD lõputööd tudengi eest kirjutada ega ümber sõnastada.
Sa annad ainult soovituslikku tagasisidet tudengi enda kirjutatud
tekstile. Lõpliku otsuse iga soovituse rakendamise kohta teeb autor.

Sinu ülesanne on tuvastada probleemid täpselt neljas kategoorias:
1. STRUKTUUR
2. AKADEEMILINE_STIIL
3. TERMINOLOOGIA
4. VIITAMISVAJADUS

[Iga kategooria kohta järgnevalt määratlus, näited ja vastunäited]

Vasta JSON-formaadis järgmise skeemi järgi:
{
  "leiud": [
    {
      "kategooria": "...",
      "tsitaat": "tekstis täpne probleemne kohaviit",
      "probleem": "lühike kirjeldus",
      "põhjendus": "miks see on probleem",
      "soovitus": "konkreetne parandusettepanek",
      "kindlus": "kõrge | keskmine | madal"
    }
  ]
}

Kui konkreetses kategoorias probleeme ei leita, jäta selle kategooria
kanded vahele. Ära leia probleeme, mida tegelikult pole.
\end{minted}

\textbf{Põhjendatud disainivalikud}:
\begin{itemize}
    \item \textbf{Eetiline meelespea}. Prompti algusesse on paigutatud selge piirang, et mudel ei tohi tudengi eest kirjutada. See on kooskõlas TÜ generatiivse tehisaru juhendiga~\cite{ut_genAI_2024} ning vähendab pilootuuringus täheldatult tõenäosust, et mudel pakub tudengile valmis ümbersõnastusi.
    \item \textbf{Tsitaadi nõue}. Iga probleemi puhul peab mudel viitama konkreetsele tekstilõigule. See raskendab \emph{halutsineerimist} ja teeb tagasiside auditeeritavaks; pilootuuringus täheldati, et tsitaadi nõue paneb mudeli oma leide täpsemini lokaliseerima.
    \item \textbf{Kindluse hinnang}. Iga leiu kohta nõutav kindlushinnang võimaldab kasutajaliideses madala kindlusega leide visuaalselt eristada ning hindamise käigus uurida, kas kindlus korreleerub leiu tegeliku õigsusega.
    \item \textbf{Põhjenduse nõue}. Mõtteahel-elemendiks olev \emph{põhjenduse} väli sunnib mudelit oma järeldust kontrollima ning aitab kasutajal otsustada, kas soovitus vastu võtta.
\end{itemize}

\subsection{Kavandatud hindamismetoodika} \label{metoodika:hindamine}

Käesolev alapeatükk kirjeldab \emph{kavandatud} täismahulise hindamise metoodikat. Selle rakendamine kogu testkogu vastu kuulub töö jätkukavasse (vt alapeatükk~\ref{hindamine:edasine}); käesoleva töö raames on metoodika kavandatud, testkogu \emph{disainitud} (vt järgmine alajaotus) ning pilootuuringu raames on kontrollitud, et prototüübi väljund on selle metoodika rakendamiseks tehniliselt sobivas vormis (JSON-struktuur, kategooriad, tsitaat, kindlushinnang).

\subsubsection{Testkogu disain}

Testkogu on kavandatud koosnema 20 eestikeelse informaatika lõputöö katkendist, kus iga katkendi pikkus on 200--400 sõna. Katkendid valitakse nii, et katavad erinevaid peatükitüüpe (vt tabel~\ref{tab:testkogu}). Allikateks on planeeritud kasutada Tartu Ülikooli arvutiteaduse instituudi avalikult kättesaadavaid bakalaureuse- ja magistritöid aastatest 2020--2024; katkendid valitakse nii, et need sisaldavad nii probleeme kui ka korrektseid lõike, ning anonüümitakse enne kuldstandardi annoteerimist. Pilootuuringus kasutati orienteeruvalt 3--5 esialgset katkendit peatükitüübi kohta, et iteratiivselt täpsustada prompti sõnastust; täismahulise testkogu lõplik koostamine ja annoteerimine on jätkukava esimene samm.

\begin{table}[htb!]
    \centering
    \caption{Kavandatud testkogu katkendite jaotus peatükitüübi järgi.}
    \label{tab:testkogu}
    \begin{tblr}{
                colspec = { Q[l] Q[c] Q[c] },
                hlines, vlines,
                row{1} = {font=\bfseries, bg = colorCellHighlight},
            }
    Peatüki tüüp & Katkendite arv & Sihtpikkus (sõnu) \\
    Sissejuhatus & 4 & 250--320 \\
    Taust ja seotud tööd & 5 & 280--360 \\
    Metoodika & 4 & 270--330 \\
    Tulemused & 4 & 240--290 \\
    Kokkuvõte & 3 & 230--260 \\
    \textbf{Kokku} & \textbf{20} & --- \\
    \end{tblr}
\end{table}

\subsubsection{Kuldstandardi koostamise protseduur}

Täismahulises hindamises annoteeritakse iga katkendi puhul probleemid kahes etapis. Esiteks märgistab autor iseseisvalt iga tuvastatud probleemi: probleemse tekstilõigu, kategooria ning lühikese põhjenduse. Teiseks vaatab teine sõltumatu annoteerija (näiteks teise õppekava tudeng või juhendaja) samad katkendid läbi, ning lahkarvamuste korral lepitakse kokku ühises lahenduses. Annoteerijatevahelist kokkulepet mõõdetakse iga kategooria kohta eraldi Coheni $\kappa$-ga~\cite{cohen_kappa_1960}, mis on kahe annoteerija ja nominaalsete kategooriate puhul konventsionaalne valik. Krippendorffi $\alpha$~\cite{krippendorff_alpha_2011} on jäetud reservi juhuks, kui hindamiskava laieneb järgmistes versioonides mitme annoteerijaga skeemiks.

\subsubsection{Mõõdikud}

Mudeli väljundi võrdlemiseks kuldstandardiga kasutatakse standardseid info\-otsingu mõõdikuid. Iga kategooria \(c\) puhul:

\begin{itemize}
    \item \textbf{Tõene positiivne} (TP\(_c\)): mudel tuvastas probleemi, mis on ka kuldstandardis.
    \item \textbf{Vale positiivne} (FP\(_c\)): mudel tuvastas probleemi, mis ei ole kuldstandardis.
    \item \textbf{Vale negatiivne} (FN\(_c\)): kuldstandardis on probleem, mida mudel ei tuvastanud.
\end{itemize}

Sellest järeldatakse järgmised mõõdikud:

\begin{equation}
    \text{Täpsus}_c = \frac{\text{TP}_c}{\text{TP}_c + \text{FP}_c}
    \label{eq:tapsus}
\end{equation}

\begin{equation}
    \text{Saagis}_c = \frac{\text{TP}_c}{\text{TP}_c + \text{FN}_c}
    \label{eq:saagis}
\end{equation}

\begin{equation}
    F_{1,c} = \frac{2 \cdot \text{Täpsus}_c \cdot \text{Saagis}_c}{\text{Täpsus}_c + \text{Saagis}_c}
    \label{eq:f1}
\end{equation}

\textbf{Kattuvuse defineerimine.} Kuna mudel ja annoteerija võivad sama probleemi kirjeldada erinevate sõnadega, ei ole võimalik nõuda eksaktset stringivastet. Kattuvust hinnatakse järgmise reegli järgi: mudeli leid loetakse tõeseks positiivseks, kui (1) tema osutatud tekstilõik kattub kuldstandardi probleemse lõiguga vähemalt 50~\% ulatuses sõnatasemel ning (2) kategooria on sama. 50~\% lävi on valitud kooskõlas span-tasemel hindamise levinud tavadega: poole ulatuses kattuv viidatud lõik tagab, et mudel ja annoteerija osutavad samale probleemile, mitte erinevatele probleemidele samas lauses. Kuna lävi ise on disainivalik, kontrollitakse selle robustsust täismahulises hindamises eraldi tundlikkusanalüüsiga, esitades peamised mõõdikud ka 30~\% ja 70~\% läviga ning jälgides, kas kvalitatiivsed järeldused muutuvad. Kui kategooria erineb, loetakse leid mudeli vaates valepositiivseks ja kuldstandardi vaates valenegatiivseks.

\subsubsection{Promptide võrdlus}

Üldist ja struktureeritud prompti võrreldakse, käivitades mudeli kõikide testkogu katkenditega kummagi prompti versioonis. Mõlemad promptid nõuavad sama JSON-skeemiga väljundit ja kategooriavalikut, mistõttu skoorimine kuldstandardi vastu toimub mõlemal juhul automaatselt, ilma käsitsi vahesammuta. Erinevus, mida hindamine mõõdab, taandub seetõttu sisuliste juhiste mõjule -- kas kategooria definitsioonid, näited ja reeglid parandavad mudeli täpsust ja saagist iga kategooria sees -- mitte väljundvormingu erinevusele. Kategooria \texttt{MUU}, mida üldine prompt lubab, kuid struktureeritud mitte, käsitatakse skoorimises automaatselt valepositiivsena, kuna see ei vasta ühelegi nelja kategooria kuldstandardi annotatsioonile.
